\documentclass[10pt, aspectratio=169]{beamer}

% --- Configuración del Tema (Estilo Atractor) ---
\usetheme{metropolis}
\usecolortheme{owl} % Fondo oscuro, alto contraste
\usepackage[spanish]{babel}
\usepackage[utf8]{inputenc}
\usepackage{amsfonts, amsmath, amssymb, amsthm}
\usepackage{graphicx}
\usepackage{tikz}
\usepackage{hyperref}
\usepackage{array}

% --- Colores Personalizados ---
\definecolor{diagonalred}{RGB}{140, 20, 40}
\definecolor{mitred}{RGB}{163, 31, 52}

% --- Información del Título ---
\title{De los Sistemas Formales a la Electrónica Aplicada}
\subtitle{Clase 2: Lógica Proposicional: Sintaxis y Semántica}
\author{TA: Rosh Guadiana P.}
\date{\today}

\begin{document}

% --- Slide 1: Título ---
\maketitle

% --- Slide 2: El Sistema MIU ---
\begin{frame}{El Sistema MIU y la Sintaxis Pura}
    \begin{columns}
        \begin{column}{0.5\textwidth}
            \textbf{Un Sistema Formal Puro}
            \begin{itemize}
                \item Los símbolos $M, I, U$ \alert{no significan nada}; son marcas vacías.
                \item La derivación opera a ciegas, separada del significado humano.
                \item \textbf{El misterio:} ¿Se puede derivar la cadena \textbf{MU}?
            \end{itemize}
            \vspace{0.2cm}
            \end{block}
        \end{column}
        \begin{column}{0.5\textwidth}
            \begin{exampleblock}{Definición del Sistema (GEB)}
                \textbf{Alfabeto:} $\{M, I, U\}$ \\
                \textbf{Axioma Inicial:} $MI$ \\
                \textbf{Reglas de Inferencia:}
                \begin{enumerate}
                    \item $xI \to xIU$
                    \item $Mx \to Mxx$
                    \item $xIIIy \to xUy$
                    \item $xUUy \to xy$
                \end{enumerate}
            \end{exampleblock}
        \end{column}
    \end{columns}
\end{frame}

% --- Slide 3: Escher y Figura/Fondo ---
\begin{frame}{Escher, Figura/Fondo y las Negaciones}
    \begin{columns}
        \begin{column}{0.55\textwidth}
            En la lógica formal existe una dualidad análoga al arte de Escher:
            \begin{itemize}
                \item \textbf{La Figura (Espacio Positivo):} El conjunto de todos los \alert{teoremas} que nuestras reglas construyen.
                \item \textbf{El Fondo (Espacio Negativo):} El conjunto de los \alert{no-teoremas} que las reglas rechazan (ej. MU).
            \end{itemize}
            \vspace{0.4cm}
            \textit{¿Acaso la figura contiene la misma información que su fondo?} \\
            \vspace{0.2cm}
            \footnotesize Generar verdades mecánicamente no nos da un método automático para identificar todas las falsedades. Esta asimetría detonará la \textbf{Incompletitud de Gödel}.
        \end{column}
        \begin{column}{0.45\textwidth}
            \centering
            % Sustituir por la imagen local de Escher cuando compiles
            \includegraphics[width=0.9\textwidth]{example-image}
            \\ \small \textit{Day and Night} - M.C. Escher
        \end{column}
    \end{columns}
\end{frame}

% --- Slide 4: Boole y la Separación ---
\begin{frame}{Boole y la Gran Separación ($\vdash$ vs $\models$)}
    \textbf{Objetivo de Boole (1854):} Matematizar las Leyes del Pensamiento. \\
    \vspace{0.3cm}
    Debemos ejecutar la \alert{Gran Separación}:
    \begin{columns}
        \begin{column}{0.45\textwidth}
            \begin{block}{Sintaxis (Derivación) $\vdash$}
                Manipulación mecánica de \textbf{Fórmulas Bien Formadas (FBF)}.
                \begin{itemize}
                    \item Reglas estructurales.
                    \item Operación puramente tipográfica.
                \end{itemize}
            \end{block}
        \end{column}
        \begin{column}{0.1\textwidth}
            \centering \Large $\Longleftrightarrow$
        \end{column}
        \begin{column}{0.45\textwidth}
            \begin{block}{Semántica (Verdad) $\models$}
                Asignación de valores del mundo real mediante \textbf{Valuaciones}.
                \begin{itemize}
                    \item Tablas de verdad.
                    \item Interpretación del significado.
                \end{itemize}
            \end{block}
        \end{column}
    \end{columns}
    \vspace{0.4cm}
    \centering
    \textbf{Teorema de Corrección:} Si manipulas correctamente los símbolos mecánicos ($\vdash$), alcanzarás una verdad semántica indudable ($\models$).
\end{frame}

% --- Slide 5: Semántica Práctica ---
\begin{frame}{Ejemplo Semántico: Modus Ponens}
    \centering
    Verificando la validez semántica del Modus Ponens ($p \to q, p \models q$) mediante tablas de verdad:
    \vspace{0.4cm}
    
    \renewcommand{\arraystretch}{1.3}
    \begin{tabular}{|c|c|c|c|c|}
        \hline
        $p$ & $q$ & $p \to q$ & $(p \to q) \land p$ & $((p \to q) \land p) \to q$ \\
        \hline
        V & V & V & V & \textbf{V} \\
        V & F & F & F & \textbf{V} \\
        F & V & V & F & \textbf{V} \\
        F & F & V & F & \textbf{V} \\
        \hline
    \end{tabular}
    \vspace{0.4cm}
    
    \textit{La tautología final confirma que la regla de inferencia sintáctica preserva la verdad semántica.}
\end{frame}

% --- Slide 6: El Clímax NAND ---
\begin{frame}{El Clímax: La Completitud Funcional del NAND}
    \begin{columns}
        \begin{column}{0.4\textwidth}
            Cualquier fórmula proposicional, por compleja que sea, puede expresarse con \textbf{un solo operador universal}: el \alert{NAND} ($p \uparrow q$). \\
            \vspace{0.4cm}
            \renewcommand{\arraystretch}{1.2}
            \begin{tabular}{|c|c|c|}
                \hline
                $p$ & $q$ & $p \uparrow q$ \\
                \hline
                V & V & F \\
                V & F & V \\
                F & V & V \\
                F & F & V \\
                \hline
            \end{tabular}
        \end{column}
        \begin{column}{0.6\textwidth}
            \begin{exampleblock}{Derivaciones Algebraicas}
                \textbf{1. Negación (NOT):} $\neg p$
                \begin{equation*}
                    \neg p \equiv p \uparrow p
                \end{equation*}
                
                \textbf{2. Conjunción (AND):} $p \land q$
                \begin{align*}
                    p \land q &\equiv \neg(p \uparrow q) \\
                              &\equiv (p \uparrow q) \uparrow (p \uparrow q)
                \end{align*}
            \end{exampleblock}
            \vspace{0.3cm}
            \footnotesize \textbf{Hacia la Sesión 8:} Abandonaremos el pizarrón. Usaremos transistores de silicio para construir físicamente esta compuerta NAND, capturando la lógica del universo dentro de la materia.
        \end{column}
    \end{columns}
\end{frame}

\end{document}
